\documentclass{article}

\usepackage{hyperref}
\usepackage{amsmath,amssymb}
\usepackage{graphicx}
\usepackage{booktabs}
\usepackage{xcolor}
\usepackage{algorithm}
\usepackage{algorithmic}
\usepackage{natbib}

\title{The Aurelian Memory Core: \\Per-Layer Differentiable Memory for Transformer Models}

\author{
  Anonymous Authors \\
  \texttt{anonymous@anonymous.com}
}

\begin{document}
\maketitle

\begin{abstract}
We present the Aurelian Memory Core (AMC), a per-layer differentiable
three-tier memory hierarchy integrated into the transformer forward pass.
Tier-1 working memory uses Mamba-2 selective state spaces; Tier-2 episodic
memory uses surprise heads and Gumbel-softmax promotion gates; Tier-3
long-term storage uses trust-aware consolidation with quarantine.
On AMC-Memory (oracle smoke, $n{=}5$), scores increase from
TODO\_BASELINE to TODO\_FULL\_AMC (see
\texttt{ablation\_scores.jsonl}); GPU-trained deltas are
\textbf{TODO\_DELTA}. Standard-benchmark parity (GSM8K, MMLU) is
\textbf{TODO\_PVALUE} until engine evaluation on trained
\textbf{TODO\_MODEL\_SIZE} checkpoints. We release reproducibility
scripts, a claims ledger, and adversarial safety probes (6/6 PASS in
structural audit).
\end{abstract}

\section{Introduction}
\label{sec:intro}

Modern language models treat memory as an afterthought: the model is
trained on one-shot next-token prediction, and retrieval is added as a
post-hoc component (RAG, MemGPT). This split creates three problems:
\begin{enumerate}
  \item \textbf{Memory is not learned}. The retrieval module cannot
        influence what the model encodes.
  \item \textbf{Memory is unsafe}. Retrieved passages are injected
        into context without trust verification---the source of many
        RAG poisoning attacks.
  \item \textbf{Memory is opaque}. There is no gradient signal
        telling the model which stored facts were helpful.
\end{enumerate}

We argue that memory should be \emph{native to the architecture}:
every layer should decide, via a learned mechanism, what to hold in
working memory, what to promote to episodic storage, and what to
consolidate into long-term store.

We introduce the \textbf{Aurelian Memory Core (AMC)}, a per-layer
differentiable three-tier memory hierarchy in which:
\begin{itemize}
  \item \textbf{Tier-1} (working memory) uses Mamba-2 selective state
        spaces~\citep{dao2024transformers} with learned decay, erase,
        and write gates.
  \item \textbf{Tier-2} (episodic memory) is gated by a learned
        surprise head and a Gumbel-softmax promotion gate.
  \item \textbf{Tier-3} (long-term store) is a trust-aware consolidated
        store with quarantine, decay, and consolidation policies.
\end{itemize}

\paragraph{Contributions.}
\begin{enumerate}
  \item A hybrid transformer with per-layer differentiable memory
        (surprise heads, promotion gates, three-tier hierarchy).
  \item Trust-aware memory contracts and constitutional alignment as
        retrieval-native safety.
  \item AMC-Memory benchmark suite and open reproducibility bundle.
\end{enumerate}

\paragraph{Outline.}
Section~\ref{sec:related} reviews related work.
Section~\ref{sec:method} describes the method.
Section~\ref{sec:training} covers training objectives.
Section~\ref{sec:experiments} presents ablations and safety audits.
Section~\ref{sec:analysis} discusses qualitative analysis.
Section~\ref{sec:conclusion} concludes.

\section{Related Work}
\label{sec:related}

\paragraph{Memory-augmented language models.}
MemGPT/Letta~\citep{packer2023memgpt}, Generative
Agents~\citep{park2023generative}, and
RAG~\citep{lewis2020rag} demonstrate multi-tier or retrieval-based
memory without differentiable promotion through the forward pass.

\paragraph{Learned memory in the forward pass.}
Titans~\citep{bulian2025titans} introduces per-layer learned memory
with surprise-driven gating. AMC extends this with explicit Tier-2/3
stores, trust contracts, and constitutional alignment.

\paragraph{Selective state space models.}
Mamba~\citep{gu2023mamba} and Mamba-2~\citep{dao2024transformers}
provide the Tier-1 substrate.

\paragraph{Constitutional AI.}
Constitutional AI~\citep{bai2022constitutional} encodes safety in reward
signals; AMC encodes principles as permanent retrievable memory.

\section{Method}
\label{sec:method}

% Evidence: src/model/amc_transformer.py; tests/model/test_amc_transformer_e2e.py

\subsection{Overview}

AMC interleaves multi-head latent attention (MLA) blocks with
per-layer selective state-space (SSM) working memory, surprise
prediction, and a differentiable promotion gate that writes to
episodic (Tier-2) and consolidated (Tier-3) stores under an SDB
propose--verify--commit contract.

\subsection{Tier-1: Per-Layer SSM Working Memory}

% Evidence: src/model/mamba2_block.py; src/model/amc_ssm_layer.py

For layer $\ell$, the SSM maintains a tensor state
$h_t^{(\ell)} \in \mathbb{R}^{H \times D \times N}$ (batch, heads,
head dim, state dim)---not a mean-pooled vector.
With learned discretization $\Delta_t$ and diagonal $A$,
zero-order-hold discretization gives
\begin{equation}
  \bar{A}_t = \exp(\Delta_t A), \qquad
  h_t = \bar{A}_t \odot h_{t-1} + x_t \otimes (\Delta_t B_t), \qquad
  y_t = C_t h_t + D \odot x_t.
\end{equation}
Gates $d_t, e_t, w_t \in [0,1]^N$ modulate decay, erase, and write
(\texttt{AMCGateController}).

\subsection{Surprise Prediction and Stop-Gradient Boundary}

% Evidence: src/model/amc_surprise.py

The surprise score uses the \emph{detached} hidden state as input:
\begin{equation}
  s_t^{(\ell)} = \sigma\!\left(f_s\!\left(\operatorname{sg}\!\left[h_t^{(\ell)}\right]\right)\right).
\end{equation}
Gradients from the surprise loss do not flow into the main block through
this head; the SSM and attention paths remain the primary representation learners.

\subsection{Differentiable Promotion Gate}

% Evidence: src/model/amc_promotion.py; tests/model/test_amc_promotion.py

Let $z_t^{\mathrm{soft}}$ be Gumbel-softmax probabilities and
$z_t^{\mathrm{hard}} \in \{0,1\}$ the hard sample.
The straight-through store estimator is
\begin{equation}
  z_t = z_t^{\mathrm{hard}} - \operatorname{sg}\!\left[z_t^{\mathrm{soft}}\right] + z_t^{\mathrm{soft}},
\end{equation}
so the forward pass uses the hard decision while gradients flow through
$z_t^{\mathrm{soft}}$.

\subsection{Tier-2 Episodic Memory and SDB Contract}

% Evidence: src/memory/sdb_runtime.py; tests/memory/test_sdb_memory_runtime.py

Episodic writes follow
\begin{equation}
  \mathrm{propose}(m) \rightarrow \mathrm{verify}(m, e) \rightarrow
  \mathrm{commit}(m) \;\text{or}\; \mathrm{reject}(m).
\end{equation}
Commit requires verification produced by \texttt{runtime.verify()};
forged verifications, payload mutation after verify, and ID mismatches
raise errors (fail-closed).

\subsection{Tier-3 Long-Term Trust-Aware Memory}

% Evidence: src/memory/amc_tier3.py; src/agent/consolidation_agent.py

Consolidation moves high-surprise episodic entries into Tier-3 with
trust level, decay policy, and provenance metadata.

\subsection{Trust-Aware Cache Identity}

% Evidence: src/serving/amc_kv_cache.py; src/memory/amc_runtime_cache.py

Prefix cache keys bind token content, tier, trust state, quarantine flag,
and revocation epoch so quarantined or revoked entries cannot alias trusted KV.

\subsection{Constitutional Memory Retrieval}

% Evidence: src/agent/constitutional_memory.py

Safety principles are permanent Tier-3 entries (non-evictable, trusted)
always eligible for retrieval during generation.

\subsection{Training Objectives (Summary)}

% Evidence: src/training/amc_losses.py

\begin{equation}
  \mathcal{L} = \alpha \mathcal{L}_{\text{SFT}} + \beta \mathcal{L}_{\text{surprise}}
  + \gamma \mathcal{L}_{\text{consistency}} + \delta \mathcal{L}_{\text{promotion}}.
\end{equation}
Separate learning rates apply to the backbone, promotion gate, and surprise head.

\section{Training}
\label{sec:training}

% Evidence: src/training/amc_trainer.py; src/training/launch_amc_training.py

We train Aurelius-Forge 1B-AMC with the composite loss in
Section~\ref{sec:method}. Optimizer groups: main model
($3\times10^{-4}$), promotion gate ($1\times10^{-4}$), surprise head
($1\times10^{-5}$). DeepSpeed ZeRO-2 is used for multi-GPU runs
(\texttt{configs/deepspeed\_zero2.json}).

Dataset packing, importance labels, and memmap loaders are documented in
the reproducibility bundle (\texttt{docs/reproducibility/README.md}).
Final checkpoint quality metrics are \textbf{TODO\_MODEL\_SIZE} until a
full GPU training run completes.

\section{Experiments}
\label{sec:experiments}

% Evidence: docs/reproducibility/results/ablation_scores.jsonl; src/eval/ablation.py

\subsection{Experimental Setup}

Four 1B-parameter configurations: \texttt{baseline} (attention-only),
\texttt{tier1\_only}, \texttt{tier12}, and \texttt{full\_amc}.
Benchmarks: AMC-Memory, RULER NIAH, LongBench-v2, GSM8K, MMLU-57
(as implemented in \texttt{src/eval/ablation.py}).

\subsection{Main Ablation (Oracle Smoke)}

Table~\ref{tab:ablation} reports scores from
\texttt{docs/reproducibility/results/ablation\_scores.jsonl}
(oracle mode, 5 seeds on AMC-Memory). GSM8K and MMLU show harness
stubs at 1.0 until \texttt{ABLATION\_MODE=engine} with trained weights.

\begin{table}[t]
\centering
\caption{Ablation on AMC-Memory (oracle smoke, $n{=}5$).
Other benchmarks: see provenance report; GSM8K/MMLU engine eval TODO.}
\label{tab:ablation}
\begin{tabular}{lccccc}
\toprule
Config & AMC-Memory & RULER & LongBench & GSM8K & MMLU \\
\midrule
Baseline    & 0.583 & TODO & TODO & TODO & TODO \\
Tier-1 only & 0.723 & TODO & TODO & TODO & TODO \\
Tier-1+2    & 0.863 & TODO & TODO & TODO & TODO \\
Full AMC    & \textbf{0.953} & TODO & TODO & TODO & TODO \\
\bottomrule
\end{tabular}
\end{table}


\subsection{Adversarial Safety}

Six probes from T29; results in
\texttt{docs/reproducibility/results/security\_audit.json}
(Table~\ref{tab:safety}).

\begin{table}[t]
\centering
\caption{Adversarial safety audit (T29, 6/6 PASS).}
\label{tab:safety}
\begin{tabular}{ll}
\toprule
Probe & Result \\
\midrule
Forged verification & Blocked \\
Mutation after verify & Blocked \\
Secret leakage in replay & Redacted \\
Memory poisoning & Quarantined / gated \\
Constitutional integrity & Preserved \\
Replay chain tamper & Detected \\
\bottomrule
\end{tabular}
\end{table}


\subsection{Reproducibility and Cross-Validation}

Local smoke cross-validation: PASS
(\texttt{docs/reproducibility/CROSS\_VALIDATION\_REPORT.md}).
External clean-machine run: \textbf{NOT EXECUTED}.

\subsection{Limitations}

Oracle ablation does not substitute for trained-checkpoint engine eval.
Long-context and standard benchmarks remain \textbf{TODO} until GPU training
and \texttt{ABLATION\_MODE=engine} complete.
P-values use paired bootstrap in \texttt{src/eval/ablation.py}, not
two-sample $t$-tests.

\section{Analysis}
\label{sec:analysis}

Qualitative analyses planned: surprise-head calibration curves,
promotion gate duty cycle per layer, and Tier-2/3 growth across sessions.
Numeric plots are \textbf{TODO} until post-training artifacts exist.

\section{Conclusion}
\label{sec:conclusion}

AMC makes memory native to the transformer via differentiable surprise,
promotion, and trust-aware consolidation, with verifiable SDB contracts
and constitutional alignment as retrieval. Future work: full GPU ablations,
public weights, and external cross-validation on a clean machine.

\appendix
\section{Reproducibility Checklist}
\label{app:repro}

See \texttt{docs/reproducibility/README.md} and
\texttt{docs/reports/AMC\_VALIDATION\_GREEN\_GATE\_20260522.md}.

\section{Claims Ledger}
\label{app:claims}

Full claim-to-evidence mapping: \texttt{paper/CLAIMS\_LEDGER.md}.


\bibliographystyle{plainnat}
\bibliography{references}

\end{document}
